\subsection{Restricciones del Proyecto DentiTech}

El desarrollo del proyecto DentiTech esta sujeto
a condiciones academicas, organizacionales,
tecnologicas y legales que deben considerarse
durante la ingenieria de requisitos.

Estas restricciones permiten establecer limites
realistas y orientar la planificacion del equipo.

\subsubsection{Restricciones de Tiempo}

El proyecto se desarrolla dentro del periodo
academico SPA 2026--2027 de la Universidad
Tecnica Estatal de Quevedo.

Las actividades deberan ajustarse al cronograma
establecido para las semanas 4 a 17, considerando
las fechas de entrega, revisiones y evaluaciones
de la asignatura Ingenieria de Requerimientos.

\subsubsection{Restricciones de Recursos Humanos}

El equipo esta integrado por cuatro estudiantes,
quienes distribuyen las actividades de analisis
del negocio, modelado del contexto, elicitacion
de requisitos, gestion e integracion documental.

La ejecucion del proyecto depende de la
coordinacion entre los integrantes y de la
disponibilidad de los representantes de la
Clinica Odontologica Medicus para participar
en entrevistas y validaciones.

\subsubsection{Restricciones Economicas}

Se priorizara una propuesta tecnologica
sostenible, evitando costos elevados de
licenciamiento y mantenimiento.

La seleccion definitiva de herramientas y
servicios dependera de los recursos disponibles
y de las necesidades confirmadas por la clinica.

No se establece un presupuesto economico
definitivo hasta disponer de informacion
validada por la organizacion cliente.

\subsubsection{Restricciones Tecnologicas}

La solucion debera considerar la infraestructura
informatica disponible en la clinica, incluyendo
equipos de computo, conectividad y mecanismos
de respaldo.

Antes de seleccionar una arquitectura definitiva,
sera necesario verificar las condiciones reales
de los equipos y el acceso a internet.

El sistema debera contemplar mecanismos de
autenticacion, control de acceso por roles y
proteccion de la informacion almacenada.

\subsubsection{Restricciones Legales y de Datos}

El tratamiento de datos personales y de salud
debera respetar la Ley Organica de Proteccion
de Datos Personales del Ecuador.

Durante las actividades academicas se evitara
incorporar nombres, diagnosticos, documentos
de identidad u otros datos sensibles de
pacientes en repositorios publicos.

Las entrevistas, observaciones y recopilacion
de evidencias deberan realizarse con las
autorizaciones correspondientes y siguiendo
el protocolo de consentimiento informado
establecido por el equipo.

\subsubsection{Restricciones Operativas}

La obtencion de informacion dependera de los
horarios y de la disponibilidad del personal
de la clinica.

Las sesiones de levantamiento de requisitos
deberan organizarse procurando no interrumpir
la atencion odontologica.

Las funcionalidades y reglas de negocio
propuestas deberan validarse antes de
considerarse requisitos definitivos.

\subsubsection{Control de Cambios del Alcance}

Las solicitudes de nuevas funcionalidades
deberan documentarse y evaluarse en funcion
de su importancia, complejidad y efecto
sobre el cronograma del proyecto.

Se dara prioridad a los procesos esenciales
de gestion odontologica, evitando incorporar
modulos no contemplados sin la aprobacion
correspondiente.
